\section{Conclusion}
\label{sec:conc}

We have presented the Leximax Universal-Regret core, a Pareto-compatible decision
rule that replaces the efficient set as the final object of multicriteria analysis
with an auditable certificate. LexUR operates as a lexicographic minimax regret selection rule: it ranks alternatives by the
lexicographically sorted vector of normalised achievement shortfalls over a declared family of
monotone probes, and returns a robust stability class (often a singleton recommendation) together with an
explanation of where it remains fragile. We proved existence and completeness,
Pareto compatibility \emph{under a separation condition}, dominated-solution
exclusion, stability of the tolerance-based rule under bounded perturbations, and
a disappointment-perturbation bound with explicit constants, and we gave a
redundancy-aware, correlation-clustering probe construction that keeps the family
linear in the number of criteria. A replicated held-out study over ten preference
families and a stochastic smart-grid case showed that LexUR attains
worst-case robustness competitive with the strongest robust
scalarisations---performing non-inferiorly to achievement scalarisation and
sampled minimax regret, while a single-point hypervolume score ($\mathrm{HV}_1$) obtains the best average rank---and does
so while reducing sensitivity to redundant criteria, certifying
its choices, and extending naturally to noisy criteria.

The central message is that robust, explainable selection can be obtained from a declared family of monotone questions rather than from elicited scalarisation weights. In structured cases (e.g.\ linear MOPs) LexUR can sometimes be computed directly, without front enumeration; in the general case, however, it remains primarily a \emph{post hoc} candidate-set selector, and we make no claim that it replaces front generation universally. This robustness also relies entirely on the reliability of the supplied normalisation bounds; if true nadirs are uncertain, the framework must be deployed conservatively using interval estimates. By relocating preference choices into a declared class of questions, LexUR turns a candidate set into a defensible,
explainable recommendation---a low-disappointment compromise across the declared
interpretations of the criteria, accompanied by a certificate that reports which
questions remain most binding.

\paragraph{Reproducibility.}
All code, generated data, and scripts that regenerate every table and figure in
this paper are released in an open-source repository (URL withheld to preserve
anonymity during review; it will be provided in the camera-ready version),
including the four-geometry benchmark, the redundancy
benchmark, the statistical tests, the sensitivity analyses, and the smart-grid
case study.
